

\chapter{Introduction}\label{ch:intro}

\section{Motivation}

In his 1977 paper, Leslie Lamport defined \emph{safety} as a property ensuring that ``something bad never happens'' \cite{Lamport1977Proving}. In systems programming, \emph{memory safety} translates to a guarantee: no out-of-bounds reads, no null-pointer dereferences, and no other memory-related \emph{undefined behaviour} (UB). In C and C++, it is the programmer's responsibility to uphold these guarantees, and failures can be severe: Heartbleed (2014) \cite{CVE-2014-0160} was an unchecked bounds error that leaked sensitive memory. This motivates \emph{executable-level memory safety} without relying only on programmer discipline. The central question of this dissertation is not only whether source programs can be checked for memory safety, but where the trust boundary should lie once those programs have been compiled.

Existing work gives two useful but incomplete answers. Source-safe languages such as Rust check strong frontend properties, but the final executable is still only as good as the compiler transformation. Verified compilers such as CompCert \cite{compcert} make the transformation trustworthy: CompCert provides a guarantee of \emph{semantic preservation} and reduces the \emph{Trusted Computing Base} (TCB) from the compiler to Rocq's kernel \cite{RocqRefman}. Since CompCert compiles C, however, semantic preservation does not by itself make unsafe source programs safe. The natural combination would be a verified compiler for a safety-enforcing source language, but this turns compiler maintenance into proof maintenance.

\begin{table}[H]
  \centering
  \footnotesize
  \renewcommand{\arraystretch}{1.02}
  \begin{tabularx}{0.96\textwidth}{@{}>{\raggedright\arraybackslash}p{0.24\textwidth}>{\raggedright\arraybackslash}p{0.34\textwidth}>{\raggedright\arraybackslash}X@{}}
    \toprule
    \textbf{Approach} & \textbf{What it gives} & \textbf{What remains exposed} \\
    \midrule
    Source-level safety & Strong frontend guarantees, as in Rust-style systems languages. & Compiler bugs can invalidate the checked guarantee. \\
    Verified compilation & A small proof-assistant kernel replaces trust in the full compiler. & Semantic preservation is not memory safety for an unsafe source language. \\
    PCC-style checking & The compiled artefact carries evidence that can be checked before execution. & Low-level code must preserve enough evidence to be checked independently. \\
    \bottomrule
  \end{tabularx}
\end{table}


\emph{Proof-carrying code} (PCC) \cite{Necula1997PCC} offers a way to move this trust boundary down to the compiled artefact without verifying the compiler itself: the producer of low-level code also ships evidence that the code obeys a safety policy, and the consumer checks this evidence before execution. One way to realise this idea is \emph{type-preserving compilation}, where type information is preserved from the source language down to assembly. A low-level type system can then encode the relevant safety properties, leaving only the verifier trusted for memory safety.

The gap addressed by this dissertation is the lack of a small, executable systems-language toolchain that combines source-level refinement checking with independently checkable low-level output, without requiring a fully verified compiler. This dissertation presents \textbf{Veritas}, which applies PCC-style checking to a refinement-typed source language. Its pipeline is:

\pipelinebox{%
\textbf{Veritas source} $\rightarrow$ typed SSA $\rightarrow$ DTAL $\rightarrow$ independent verifier $\rightarrow$ executable
}

The target language is a version of Dependently Typed Assembly Language (DTAL), inspired by Xi and Harper \cite{Xi2001DTAL}. The novelty is in the combination: Veritas lowers source-level refinement guarantees through typed SSA into physical DTAL, including registers, spills, calling conventions, ownership operations, and ELF output, then re-checks the low-level artefact before execution. The central claim is not that the compiler is correct, but that its output can be checked independently, reducing the TCB to the standalone verifier and its dependencies.


\section{Approach and aims}

Veritas consists of four core components.

\componentblock{The Veritas Language}{A systems language with refinement types, loop invariants, and function contracts. These annotations state the source-level safety facts that must survive compilation.}
\componentblock{Dependently Typed Assembly Language}{A target language based on Xi and Harper's DTAL \cite{Xi2001DTAL} that embeds type states and constraints into the instruction stream, making the relevant safety evidence explicit.}
\componentblock{Type-preserving compiler pipeline}{A pipeline that lowers source programs through typed SSA (Static Single Assignment) to DTAL while preserving safety invariants. Its output is accepted only if the verifier can check it.}
\componentblock{DTAL verifier}{A tool that checks DTAL programs without access to the source code. This is the main trust-boundary shift: the final safety decision is made by the verifier rather than by the compiler pipeline.}

\noindent The trust argument above gives two axes for comparing toolchains. Figure~\ref{fig:safety-sem-taxonomy} separates two guarantees that are often conflated: semantic preservation and executable-level memory safety.

\begin{figure}[H]
  \centering
  \resizebox{0.62\textwidth}{!}{
    \begin{tikzpicture}[
        box/.style={draw, thick, minimum width=4.5cm, minimum height=2.5cm, align=center, font=\sffamily},
        header/.style={font=\sffamily\bfseries, align=center}
      ]

      \node[box] (c1r1) {C / C++ \\ Standard Rust};
      \node[box, right=0.12cm of c1r1] (c2r1) {Veritas \\ (PCC / DTAL)};
      \node[box, below=0.12cm of c1r1] (c1r2) {CompCert};
      \node[box, right=0.12cm of c1r2] (c2r2) {RustCompCert\footnotemark};

      \node[header, above=0.3cm of c1r1] {Not Memory Safe\\(Executable Level)};
      \node[header, above=0.3cm of c2r1] {Memory Safe\\(Executable Level)};

      \node[header, left=0.3cm of c1r1, align=right] {No Semantic\\Preservation};
      \node[header, left=0.3cm of c1r2, align=right] {Semantic\\Preservation};
    \end{tikzpicture}
  }
  \caption{A taxonomy of systems programming toolchains based on their executable-level guarantees.}
  \label{fig:safety-sem-taxonomy}
\end{figure}
\footnotetext{RustCompCert \cite{RustCompCert} is an ongoing project to verify an end-to-end compiler for the Rust programming language, but it currently only focuses on a small sequential subset of the language.}


Together, these components support the core question of the dissertation: whether a small systems-language toolchain can realise the trust model suggested by proof-carrying code, with an untrusted compiler whose low-level output is independently checked before execution. The dissertation evaluates this along two main axes:

\componentblock{Trust architecture}{Can the relevant low-level safety properties of compiled programs be checked independently at the DTAL level, and does the verifier remain small enough to audit without access to the source code or compiler state?}
\componentblock{Expressiveness and practicality}{Can the source language support realistic imperative, systems-oriented programs while retaining automatic verification?}
